% --- Archon physics formalization source begin ---
% archon:physics
% archon:covers PhyXMiniProblems/problem_phyx_mini_0975.lean
% archon:source-report reports/phyx_mini/problem_phyx_mini_0975.source.json

\chapter{Physics problem phyx\_mini\_0975}
\label{ch:PhyXMiniProblems_problem_phyx_mini_0975}

\paragraph{Problem source.}
\#\# Physical scenario

A very long, cylindrical wire of radius \textbackslash{}( R \textbackslash{}) carries a current \textbackslash{}( I\_0 \textbackslash{}) uniformly distributed across the cross section of the wire.

\#\# Question

Calculate the magnetic flux through a rectangle that has one side of length \textbackslash{}( W \textbackslash{}) running down the center of the wire and another side of length \textbackslash{}( R \textbackslash{}).

\#\# Answer choices

- A: A: \textbackslash{}frac\{\{\textbackslash{}mu\_0 I\_0 W\}\}\{\{2\textbackslash{}pi R\}\}

- B: B: \textbackslash{}frac\{\{\textbackslash{}mu\_0 I W\}\}\{\{4\textbackslash{}pi\}\}

- C: C: \textbackslash{}frac\{\{\textbackslash{}mu\_0 I\_0 R W\}\}\{\{2\textbackslash{}pi\}\}

- D: D: \textbackslash{}frac\{\{\textbackslash{}mu\_0 I\_0 W\}\}\{\{2\textbackslash{}pi\}\}

\#\# Figure caption (auxiliary; use the image as primary evidence)

The image depicts a three-dimensional representation of a cylindrical object. Here's a breakdown of its components and relationships:

1. **Cylinder**: 
   - The primary shape is a transparent cylinder oriented horizontally.
   - The cylinder has an open or partially open structure, showing its inside.

2. **Rectangular Strip**: 
   - Inside the cylinder, there's a red rectangular strip that runs longitudinally along its interior.
   - The strip is centrally aligned within the cylinder.

3. **Labels and Measurements**:
   - "I₀": This label is placed near the circular end of the cylinder.
   - "R": Denotes the radial distance from the center to the edge of the circular base, marked with a double-headed arrow.
   - "W": Represents the width along the length of the cylinder, indicated by a double-headed arrow spanning the inside lengthwise.

The illustration likely serves to represent a cross-sectional view of a physical concept or component, with dimensions and components highlighted for educational or analytical purposes.

\#\# Dataset metadata

Magnetism; Physical Model Grounding Reasoning, Numerical Reasoning

\paragraph{Recorded answer/context.}
B: B: \textbackslash{}frac\{\{\textbackslash{}mu\_0 I W\}\}\{\{4\textbackslash{}pi\}\}

\paragraph{Figure/image path.}
/root/proposal\_for\_physic/science-mango/phyx\_mini\_run/phyx\_data/test\_image/975.png

\paragraph{Formalization target.}
create a compiling Lean file with sorry bodies at `PhyXMiniProblems/problem\_phyx\_mini\_0975.lean`. The Lean declarations must preserve the physical quantities, dimensions or dimensional roles, figure labels, governing-law hypotheses, and final relation expressed by this problem.
Use Mathlib/Physlib names found through LeanExplore where available. If a physics API is missing, introduce faithful local abstractions rather than scalar placeholder aliases.

\begin{theorem}[Physics formalization target]
\label{thm:physics:phyx_mini_0975:target}
\lean{PhyXMiniProblems.ProblemPhyXMini0975.problem_phyx_mini_0975}
\uses{def:physics:phyx-mini-0975:phyxminiproblems-problemphyxmini0975-lengthinmeters, def:physics:phyx-mini-0975:phyxminiproblems-problemphyxmini0975-currentinamperes, def:physics:phyx-mini-0975:phyxminiproblems-problemphyxmini0975-magneticfluxinwebers, def:physics:phyx-mini-0975:phyxminiproblems-problemphyxmini0975-permeabilityinhenriespermeter, def:physics:phyx-mini-0975:phyxminiproblems-problemphyxmini0975-infinitecylindricalwirefluxsetup, def:physics:phyx-mini-0975:phyxminiproblems-problemphyxmini0975-matchesproblemandprimaryfigure, def:physics:phyx-mini-0975:phyxminiproblems-problemphyxmini0975-hasphysicalinfinitewireconfiguration, def:physics:phyx-mini-0975:phyxminiproblems-problemphyxmini0975-satisfiesexactinfinitecylindricalgeometry, def:physics:phyx-mini-0975:phyxminiproblems-problemphyxmini0975-satisfiesradiallongitudinalrectanglegeometry, def:physics:phyx-mini-0975:phyxminiproblems-problemphyxmini0975-satisfiesuniformcrosssectioncurrentlaws, def:physics:phyx-mini-0975:phyxminiproblems-problemphyxmini0975-satisfiesinteriorinfinitewireamperelaws, def:physics:phyx-mini-0975:phyxminiproblems-problemphyxmini0975-satisfiesradialrectanglemagneticfluxlaw}
The assigned autoformalize agent should translate this physics problem into Lean declarations in the covered file, with theorem and lemma proof bodies written as `by sorry`.
\end{theorem}
\begin{proof}
This is an autoformalization task, not a proof task. Produce statements that can later be proved by the physics prover without weakening the physical model.
\end{proof}

% --- Archon declaration topology begin ---
\section{Declaration topology}
\begin{definition}[electric Current Dimension]
  \label{def:physics:phyx-mini-0975:phyxminiproblems-problemphyxmini0975-electriccurrentdimension}
  \lean{PhyXMiniProblems.ProblemPhyXMini0975.electricCurrentDimension}
  Electric current has dimension C T⁻¹ (ampere)
\end{definition}
\begin{proof}
  This is the declared model definition of electric Current Dimension.
\end{proof}
\begin{definition}[current Density Dimension]
  \label{def:physics:phyx-mini-0975:phyxminiproblems-problemphyxmini0975-currentdensitydimension}
  \lean{PhyXMiniProblems.ProblemPhyXMini0975.currentDensityDimension}
  \uses{def:physics:phyx-mini-0975:phyxminiproblems-problemphyxmini0975-electriccurrentdimension}
  Current density has dimension electric current per area
\end{definition}
\begin{proof}
  This is the declared model definition of current Density Dimension.
\end{proof}
\begin{definition}[magnetic Flux Density Dimension]
  \label{def:physics:phyx-mini-0975:phyxminiproblems-problemphyxmini0975-magneticfluxdensitydimension}
  \lean{PhyXMiniProblems.ProblemPhyXMini0975.magneticFluxDensityDimension}
  Magnetic flux density has dimension M T⁻¹ C⁻¹ (tesla)
\end{definition}
\begin{proof}
  This is the declared model definition of magnetic Flux Density Dimension.
\end{proof}
\begin{definition}[magnetic Flux Dimension]
  \label{def:physics:phyx-mini-0975:phyxminiproblems-problemphyxmini0975-magneticfluxdimension}
  \lean{PhyXMiniProblems.ProblemPhyXMini0975.magneticFluxDimension}
  \uses{def:physics:phyx-mini-0975:phyxminiproblems-problemphyxmini0975-magneticfluxdensitydimension}
  Magnetic flux has dimension magnetic flux density times area (weber)
\end{definition}
\begin{proof}
  This is the declared model definition of magnetic Flux Dimension.
\end{proof}
\begin{definition}[magnetic Permeability Dimension]
  \label{def:physics:phyx-mini-0975:phyxminiproblems-problemphyxmini0975-magneticpermeabilitydimension}
  \lean{PhyXMiniProblems.ProblemPhyXMini0975.magneticPermeabilityDimension}
  Magnetic permeability has dimension M L C⁻² (henry per metre)
\end{definition}
\begin{proof}
  This is the declared model definition of magnetic Permeability Dimension.
\end{proof}
\begin{definition}[Length Quantity]
  \label{def:physics:phyx-mini-0975:phyxminiproblems-problemphyxmini0975-lengthquantity}
  \lean{PhyXMiniProblems.ProblemPhyXMini0975.LengthQuantity}
  A nonnegative, unit-independent physical length
\end{definition}
\begin{proof}
  This is the declared model definition of Length Quantity.
\end{proof}
\begin{definition}[Electric Current Magnitude]
  \label{def:physics:phyx-mini-0975:phyxminiproblems-problemphyxmini0975-electriccurrentmagnitude}
  \lean{PhyXMiniProblems.ProblemPhyXMini0975.ElectricCurrentMagnitude}
  \uses{def:physics:phyx-mini-0975:phyxminiproblems-problemphyxmini0975-electriccurrentdimension}
  A nonnegative, unit-independent electric-current magnitude
\end{definition}
\begin{proof}
  This is the declared model definition of Electric Current Magnitude.
\end{proof}
\begin{definition}[Current Density Magnitude]
  \label{def:physics:phyx-mini-0975:phyxminiproblems-problemphyxmini0975-currentdensitymagnitude}
  \lean{PhyXMiniProblems.ProblemPhyXMini0975.CurrentDensityMagnitude}
  \uses{def:physics:phyx-mini-0975:phyxminiproblems-problemphyxmini0975-currentdensitydimension}
  A nonnegative, unit-independent current-density magnitude
\end{definition}
\begin{proof}
  This is the declared model definition of Current Density Magnitude.
\end{proof}
\begin{definition}[Magnetic Flux Density Magnitude]
  \label{def:physics:phyx-mini-0975:phyxminiproblems-problemphyxmini0975-magneticfluxdensitymagnitude}
  \lean{PhyXMiniProblems.ProblemPhyXMini0975.MagneticFluxDensityMagnitude}
  \uses{def:physics:phyx-mini-0975:phyxminiproblems-problemphyxmini0975-magneticfluxdensitydimension}
  A nonnegative, unit-independent magnetic-flux-density magnitude
\end{definition}
\begin{proof}
  This is the declared model definition of Magnetic Flux Density Magnitude.
\end{proof}
\begin{definition}[Magnetic Flux Magnitude]
  \label{def:physics:phyx-mini-0975:phyxminiproblems-problemphyxmini0975-magneticfluxmagnitude}
  \lean{PhyXMiniProblems.ProblemPhyXMini0975.MagneticFluxMagnitude}
  \uses{def:physics:phyx-mini-0975:phyxminiproblems-problemphyxmini0975-magneticfluxdimension}
  A nonnegative magnetic-flux magnitude through the displayed rectangle
\end{definition}
\begin{proof}
  This is the declared model definition of Magnetic Flux Magnitude.
\end{proof}
\begin{definition}[Magnetic Permeability Magnitude]
  \label{def:physics:phyx-mini-0975:phyxminiproblems-problemphyxmini0975-magneticpermeabilitymagnitude}
  \lean{PhyXMiniProblems.ProblemPhyXMini0975.MagneticPermeabilityMagnitude}
  \uses{def:physics:phyx-mini-0975:phyxminiproblems-problemphyxmini0975-magneticpermeabilitydimension}
  A nonnegative, unit-independent magnetic permeability
\end{definition}
\begin{proof}
  This is the declared model definition of Magnetic Permeability Magnitude.
\end{proof}
\begin{definition}[nonnegative SIReadout]
  \label{def:physics:phyx-mini-0975:phyxminiproblems-problemphyxmini0975-nonnegativesireadout}
  \lean{PhyXMiniProblems.ProblemPhyXMini0975.nonnegativeSIReadout}
  Read a nonnegative dimensionful quantity in coherent SI units
\end{definition}
\begin{proof}
  This is the declared model definition of nonnegative SIReadout.
\end{proof}
\begin{definition}[length In Meters]
  \label{def:physics:phyx-mini-0975:phyxminiproblems-problemphyxmini0975-lengthinmeters}
  \lean{PhyXMiniProblems.ProblemPhyXMini0975.lengthInMeters}
  \uses{def:physics:phyx-mini-0975:phyxminiproblems-problemphyxmini0975-lengthquantity, def:physics:phyx-mini-0975:phyxminiproblems-problemphyxmini0975-nonnegativesireadout}
  Read a length in metres
\end{definition}
\begin{proof}
  This is the declared model definition of length In Meters.
\end{proof}
\begin{definition}[current In Amperes]
  \label{def:physics:phyx-mini-0975:phyxminiproblems-problemphyxmini0975-currentinamperes}
  \lean{PhyXMiniProblems.ProblemPhyXMini0975.currentInAmperes}
  \uses{def:physics:phyx-mini-0975:phyxminiproblems-problemphyxmini0975-electriccurrentmagnitude, def:physics:phyx-mini-0975:phyxminiproblems-problemphyxmini0975-nonnegativesireadout}
  Read an electric-current magnitude in amperes
\end{definition}
\begin{proof}
  This is the declared model definition of current In Amperes.
\end{proof}
\begin{definition}[current Density In Amperes Per Square Meter]
  \label{def:physics:phyx-mini-0975:phyxminiproblems-problemphyxmini0975-currentdensityinamperespersquaremeter}
  \lean{PhyXMiniProblems.ProblemPhyXMini0975.currentDensityInAmperesPerSquareMeter}
  \uses{def:physics:phyx-mini-0975:phyxminiproblems-problemphyxmini0975-currentdensitymagnitude, def:physics:phyx-mini-0975:phyxminiproblems-problemphyxmini0975-nonnegativesireadout}
  Read a current-density magnitude in amperes per square metre
\end{definition}
\begin{proof}
  This is the declared model definition of current Density In Amperes Per Square Meter.
\end{proof}
\begin{definition}[magnetic Flux Density In Teslas]
  \label{def:physics:phyx-mini-0975:phyxminiproblems-problemphyxmini0975-magneticfluxdensityinteslas}
  \lean{PhyXMiniProblems.ProblemPhyXMini0975.magneticFluxDensityInTeslas}
  \uses{def:physics:phyx-mini-0975:phyxminiproblems-problemphyxmini0975-magneticfluxdensitymagnitude, def:physics:phyx-mini-0975:phyxminiproblems-problemphyxmini0975-nonnegativesireadout}
  Read a magnetic-flux-density magnitude in teslas
\end{definition}
\begin{proof}
  This is the declared model definition of magnetic Flux Density In Teslas.
\end{proof}
\begin{definition}[magnetic Flux In Webers]
  \label{def:physics:phyx-mini-0975:phyxminiproblems-problemphyxmini0975-magneticfluxinwebers}
  \lean{PhyXMiniProblems.ProblemPhyXMini0975.magneticFluxInWebers}
  \uses{def:physics:phyx-mini-0975:phyxminiproblems-problemphyxmini0975-magneticfluxmagnitude, def:physics:phyx-mini-0975:phyxminiproblems-problemphyxmini0975-nonnegativesireadout}
  Read magnetic flux in webers
\end{definition}
\begin{proof}
  This is the declared model definition of magnetic Flux In Webers.
\end{proof}
\begin{definition}[permeability In Henries Per Meter]
  \label{def:physics:phyx-mini-0975:phyxminiproblems-problemphyxmini0975-permeabilityinhenriespermeter}
  \lean{PhyXMiniProblems.ProblemPhyXMini0975.permeabilityInHenriesPerMeter}
  \uses{def:physics:phyx-mini-0975:phyxminiproblems-problemphyxmini0975-magneticpermeabilitymagnitude, def:physics:phyx-mini-0975:phyxminiproblems-problemphyxmini0975-nonnegativesireadout}
  Read magnetic permeability in henries per metre
\end{definition}
\begin{proof}
  This is the declared model definition of permeability In Henries Per Meter.
\end{proof}
\begin{definition}[Figure Feature]
  \label{def:physics:phyx-mini-0975:phyxminiproblems-problemphyxmini0975-figurefeature}
  \lean{PhyXMiniProblems.ProblemPhyXMini0975.FigureFeature}
  Named physical features visible in phyx\_data/test\_image/975.png
\end{definition}
\begin{proof}
  This is the declared model definition of Figure Feature.
\end{proof}
\begin{definition}[Figure Label]
  \label{def:physics:phyx-mini-0975:phyxminiproblems-problemphyxmini0975-figurelabel}
  \lean{PhyXMiniProblems.ProblemPhyXMini0975.FigureLabel}
  The three labels printed in the source image
\end{definition}
\begin{proof}
  This is the declared model definition of Figure Label.
\end{proof}
\begin{definition}[Figure Label Role]
  \label{def:physics:phyx-mini-0975:phyxminiproblems-problemphyxmini0975-figurelabelrole}
  \lean{PhyXMiniProblems.ProblemPhyXMini0975.FigureLabelRole}
  Physical role of each printed label
\end{definition}
\begin{proof}
  This is the declared model definition of Figure Label Role.
\end{proof}
\begin{definition}[Cylindrical Wire Figure]
  \label{def:physics:phyx-mini-0975:phyxminiproblems-problemphyxmini0975-cylindricalwirefigure}
  \lean{PhyXMiniProblems.ProblemPhyXMini0975.CylindricalWireFigure}
  \uses{def:physics:phyx-mini-0975:phyxminiproblems-problemphyxmini0975-figurefeature, def:physics:phyx-mini-0975:phyxminiproblems-problemphyxmini0975-figurelabel, def:physics:phyx-mini-0975:phyxminiproblems-problemphyxmini0975-figurelabelrole}
  Literal and qualitative evidence transcribed from the primary image
\end{definition}
\begin{proof}
  This is the declared model definition of Cylindrical Wire Figure.
\end{proof}
\begin{definition}[Radial Longitudinal Rectangle]
  \label{def:physics:phyx-mini-0975:phyxminiproblems-problemphyxmini0975-radiallongitudinalrectangle}
  \lean{PhyXMiniProblems.ProblemPhyXMini0975.RadialLongitudinalRectangle}
  \uses{def:physics:phyx-mini-0975:phyxminiproblems-problemphyxmini0975-lengthquantity}
  The ambient radial-longitudinal half-plane is parametrized by radial distance and axial distance in metres. Restricting the first coordinate to [0, R] and the second to [0, W] gives the red spanning surface. Keeping the map defined for every axial coordinate also lets the infinite-wire field law state exact axial translation invariance
\end{definition}
\begin{proof}
  This is the declared model definition of Radial Longitudinal Rectangle.
\end{proof}
\begin{definition}[Infinite Cylindrical Wire Flux Setup]
  \label{def:physics:phyx-mini-0975:phyxminiproblems-problemphyxmini0975-infinitecylindricalwirefluxsetup}
  \lean{PhyXMiniProblems.ProblemPhyXMini0975.InfiniteCylindricalWireFluxSetup}
  \uses{def:physics:phyx-mini-0975:phyxminiproblems-problemphyxmini0975-lengthquantity, def:physics:phyx-mini-0975:phyxminiproblems-problemphyxmini0975-electriccurrentmagnitude, def:physics:phyx-mini-0975:phyxminiproblems-problemphyxmini0975-currentdensitymagnitude, def:physics:phyx-mini-0975:phyxminiproblems-problemphyxmini0975-magneticfluxdensitymagnitude, def:physics:phyx-mini-0975:phyxminiproblems-problemphyxmini0975-magneticfluxmagnitude, def:physics:phyx-mini-0975:phyxminiproblems-problemphyxmini0975-magneticpermeabilitymagnitude, def:physics:phyx-mini-0975:phyxminiproblems-problemphyxmini0975-cylindricalwirefigure, def:physics:phyx-mini-0975:phyxminiproblems-problemphyxmini0975-radiallongitudinalrectangle}
  Independent physical data and observables. The scalar radial profiles are dimensionful magnitudes used to calibrate Physlib's vector fields in SI units. Neither rectangleFlux nor either field profile is defined by the requested closed-form answer
\end{definition}
\begin{proof}
  This is the declared model definition of Infinite Cylindrical Wire Flux Setup.
\end{proof}
\begin{definition}[Matches Problem And Primary Figure]
  \label{def:physics:phyx-mini-0975:phyxminiproblems-problemphyxmini0975-matchesproblemandprimaryfigure}
  \lean{PhyXMiniProblems.ProblemPhyXMini0975.MatchesProblemAndPrimaryFigure}
  \uses{def:physics:phyx-mini-0975:phyxminiproblems-problemphyxmini0975-figurefeature, def:physics:phyx-mini-0975:phyxminiproblems-problemphyxmini0975-figurelabel, def:physics:phyx-mini-0975:phyxminiproblems-problemphyxmini0975-infinitecylindricalwirefluxsetup}
  The prose supplies the uniform-current scenario; the bitmap supplies the named features and label roles. The phrase "very long" is handled separately by SatisfiesExactInfiniteCylindricalGeometry, rather than by a qualitative regime field in this readout structure. The side marked R is identified with the physical wire radius, and the side marked W is the rectangle's axial length. No magnetic-flux result is stored here
\end{definition}
\begin{proof}
  This is the declared model definition of Matches Problem And Primary Figure.
\end{proof}
\begin{definition}[Has Physical Infinite Wire Configuration]
  \label{def:physics:phyx-mini-0975:phyxminiproblems-problemphyxmini0975-hasphysicalinfinitewireconfiguration}
  \lean{PhyXMiniProblems.ProblemPhyXMini0975.HasPhysicalInfiniteWireConfiguration}
  \uses{def:physics:phyx-mini-0975:phyxminiproblems-problemphyxmini0975-lengthinmeters, def:physics:phyx-mini-0975:phyxminiproblems-problemphyxmini0975-currentinamperes, def:physics:phyx-mini-0975:phyxminiproblems-problemphyxmini0975-currentdensityinamperespersquaremeter, def:physics:phyx-mini-0975:phyxminiproblems-problemphyxmini0975-permeabilityinhenriespermeter, def:physics:phyx-mini-0975:phyxminiproblems-problemphyxmini0975-infinitecylindricalwirefluxsetup}
  Positivity and non-vacuity conditions for the physical configuration
\end{definition}
\begin{proof}
  This is the declared model definition of Has Physical Infinite Wire Configuration.
\end{proof}
\begin{definition}[Satisfies Exact Infinite Cylindrical Geometry]
  \label{def:physics:phyx-mini-0975:phyxminiproblems-problemphyxmini0975-satisfiesexactinfinitecylindricalgeometry}
  \lean{PhyXMiniProblems.ProblemPhyXMini0975.SatisfiesExactInfiniteCylindricalGeometry}
  \uses{def:physics:phyx-mini-0975:phyxminiproblems-problemphyxmini0975-lengthinmeters, def:physics:phyx-mini-0975:phyxminiproblems-problemphyxmini0975-infinitecylindricalwirefluxsetup}
  Exact geometric contract replacing the source's informal "very long" phrase. The axis contains a point at every real axial coordinate, and the solid cylinder has no axial cutoff: interior and boundary membership depend only on distance from the axis. This is an exact infinite-cylinder model, not a finite-length regime flag used to assert an approximate field as an equality
\end{definition}
\begin{proof}
  This is the declared model definition of Satisfies Exact Infinite Cylindrical Geometry.
\end{proof}
\begin{definition}[Satisfies Radial Longitudinal Rectangle Geometry]
  \label{def:physics:phyx-mini-0975:phyxminiproblems-problemphyxmini0975-satisfiesradiallongitudinalrectanglegeometry}
  \lean{PhyXMiniProblems.ProblemPhyXMini0975.SatisfiesRadialLongitudinalRectangleGeometry}
  \uses{def:physics:phyx-mini-0975:phyxminiproblems-problemphyxmini0975-lengthinmeters, def:physics:phyx-mini-0975:phyxminiproblems-problemphyxmini0975-infinitecylindricalwirefluxsetup}
  Geometric meaning of the centered radial-longitudinal rectangle. The edge at radial coordinate zero lies on the wire axis; the opposite edge lies on the cylindrical boundary at radius R. Its oriented normal is perpendicular to both in-surface directions
\end{definition}
\begin{proof}
  This is the declared model definition of Satisfies Radial Longitudinal Rectangle Geometry.
\end{proof}
\begin{definition}[Satisfies Uniform Cross Section Current Laws]
  \label{def:physics:phyx-mini-0975:phyxminiproblems-problemphyxmini0975-satisfiesuniformcrosssectioncurrentlaws}
  \lean{PhyXMiniProblems.ProblemPhyXMini0975.SatisfiesUniformCrossSectionCurrentLaws}
  \uses{def:physics:phyx-mini-0975:phyxminiproblems-problemphyxmini0975-lengthinmeters, def:physics:phyx-mini-0975:phyxminiproblems-problemphyxmini0975-currentinamperes, def:physics:phyx-mini-0975:phyxminiproblems-problemphyxmini0975-currentdensityinamperespersquaremeter, def:physics:phyx-mini-0975:phyxminiproblems-problemphyxmini0975-infinitecylindricalwirefluxsetup}
  Uniform axial current density inside the solid wire and its circular-area integrals. The total-current relation uses the full disk of radius R; the enclosed-current relation uses a concentric disk of arbitrary interior radius. These are general current-distribution laws and mention no magnetic flux
\end{definition}
\begin{proof}
  This is the declared model definition of Satisfies Uniform Cross Section Current Laws.
\end{proof}
\begin{definition}[Satisfies Interior Infinite Wire Ampere Laws]
  \label{def:physics:phyx-mini-0975:phyxminiproblems-problemphyxmini0975-satisfiesinteriorinfinitewireamperelaws}
  \lean{PhyXMiniProblems.ProblemPhyXMini0975.SatisfiesInteriorInfiniteWireAmpereLaws}
  \uses{def:physics:phyx-mini-0975:phyxminiproblems-problemphyxmini0975-lengthinmeters, def:physics:phyx-mini-0975:phyxminiproblems-problemphyxmini0975-currentinamperes, def:physics:phyx-mini-0975:phyxminiproblems-problemphyxmini0975-magneticfluxdensityinteslas, def:physics:phyx-mini-0975:phyxminiproblems-problemphyxmini0975-permeabilityinhenriespermeter, def:physics:phyx-mini-0975:phyxminiproblems-problemphyxmini0975-infinitecylindricalwirefluxsetup}
  Cylindrical symmetry and Ampere's law inside the wire. On the selected radial-longitudinal half-plane, the azimuthal magnetic field is parallel to the rectangle's oriented normal and is independent of axial position. At positive radius the circular Amperian loop has circumference 2 pi r
\end{definition}
\begin{proof}
  This is the declared model definition of Satisfies Interior Infinite Wire Ampere Laws.
\end{proof}
\begin{definition}[Satisfies Radial Rectangle Magnetic Flux Law]
  \label{def:physics:phyx-mini-0975:phyxminiproblems-problemphyxmini0975-satisfiesradialrectanglemagneticfluxlaw}
  \lean{PhyXMiniProblems.ProblemPhyXMini0975.SatisfiesRadialRectangleMagneticFluxLaw}
  \uses{def:physics:phyx-mini-0975:phyxminiproblems-problemphyxmini0975-lengthinmeters, def:physics:phyx-mini-0975:phyxminiproblems-problemphyxmini0975-magneticfluxdensityinteslas, def:physics:phyx-mini-0975:phyxminiproblems-problemphyxmini0975-magneticfluxinwebers, def:physics:phyx-mini-0975:phyxminiproblems-problemphyxmini0975-infinitecylindricalwirefluxsetup}
  For the radial-longitudinal surface, dA = d radial * d axial. Axial translation invariance makes the surface integral equal to W times the radial integral of the aligned azimuthal field magnitude. This is the general definition/law for the observable flux, not its requested evaluated value
\end{definition}
\begin{proof}
  This is the declared model definition of Satisfies Radial Rectangle Magnetic Flux Law.
\end{proof}
% --- Archon declaration topology end ---
% --- Archon physics formalization source end ---
